% Results: the execution evidence and the figures.

This section reports the executed evidence in five subsections that follow the
execution record, then a sixth on the per commit processing evidence. Every
number is reported exactly as it appears in the execution files, which are mapped
in \autoref{tab:results_sources}. The platform was executed autonomously by
Claude Code Opus 4.7 (1M context) Max, and across the whole run fifteen example
scripts ran to a zero exit code, the test suite passed, the static analysis was
clean, and the gate and factory safety surfaces were exercised in
full~\cite{repo-cancer-automated}.

\begin{table}[H]
\caption{Results source map (abbreviated paths under the repository root).}
\label{tab:results_sources}
\centering
\begin{tabularx}{\textwidth}{L{4.7cm} >{\raggedright\arraybackslash}X}
\toprule
\textbf{Source (under repo root)} & \textbf{Use for} \\
\midrule
  \texttt{execution/01-foundation/} : \texttt{test-suite.md} & 51 of 51 tests, the 8 module counts, the verbatim roster. \\
  \texttt{execution/02-pipeline/} : README and 3 artifacts & Five deliverables, the 2.5x factor, the three generated artifacts. \\
  \texttt{execution/03-vvuq/} : \texttt{README.md} & Examples 1 to 3 and the full 6 case gate decision surface. \\
  \texttt{execution/04-stage1-automation/} : README and chunk artifact & Simulation, ingestion, chunking, and scheduler outcomes, and the multibyte finding. \\
  \texttt{execution/05-physical-ai-stage2/} : README and PDAC timeline & The PDAC pilot timeline and the lights off safety surface. \\
\bottomrule
\end{tabularx}
\end{table}

\subsection{Foundation and the verification floor}

Before any deliverable was generated, the verification floor was established. The
full automated test suite ran on the container Python, version 3.11.15, under
pytest 9.0.3, and all 51 collected tests passed with none failing and none
skipped. The suite was also run with the optional numerical backend uninstalled,
and the result was identical, which confirms that no test silently depends on an
optional package. The 51 tests are distributed across eight modules that map one
to one onto the source packages, as shown in \autoref{tab:results_tests} and
visualized in \autoref{figtreemap}. Alongside the tests, the three static
analysis checks, the lint check, the format check, and the YAML lint, were all
clean. This is the verification half of VVUQ applied to the platform itself, and
it held without exception, which is the precondition for trusting anything the
platform then produces.

\begin{table}[H]
\caption{Automated tests by module (51 of 51 passed, 0 skipped).}
\label{tab:results_tests}
\centering
\begin{tabularx}{\textwidth}{L{4.2cm} C{1.8cm} >{\raggedright\arraybackslash}X}
\toprule
\textbf{Test module} & \textbf{Tests} & \textbf{Target package} \\
\midrule
  \texttt{test\_foundation.py}   & 17 & repository files, configs, verify script \\
  \texttt{test\_vvuq.py}         & 8  & verification, validation, uncertainty, gate \\
  \texttt{test\_pipeline.py}     & 5  & pipeline, the five established methods \\
  \texttt{test\_physical\_ai.py} & 5  & lights off factory, hybrid pilot \\
  \texttt{test\_simulation.py}   & 4  & triple runner, consensus \\
  \texttt{test\_ingestion.py}    & 4  & web search, PDF processor \\
  \texttt{test\_chunking.py}     & 4  & chunker, README generator \\
  \texttt{test\_scheduler.py}    & 4  & commit scheduler \\
  \midrule
  Total                          & 51 & all eight packages \\
\bottomrule
\end{tabularx}
\end{table}

\begin{figure}[H]
    \hspace{-0.8cm}%
    \IfFileExists{Images/test-coverage-treemap.png}
      {\includegraphics[width=1.1\linewidth, trim={0 0 0 0}, clip]{Images/test-coverage-treemap.png}}
      {\figslot{Test coverage treemap, 51 tests across 8 modules}}
    \caption{Test coverage treemap of the 51 passing tests across the eight modules, each tile sized to its test count.}
    \label{figtreemap}
\end{figure}

\subsection{Pipeline generation}

With the floor established, the pipeline produced its deliverables. Across the
three pipeline examples, five deliverables completed all four producing stages,
and every deliverable reported the same acceleration factor of 2.5, computed from
the audited formula at the default of three simulation runs, that is 30 baseline
days reduced to 12 automated days. This decomposition is shown in
\autoref{figwaterfall}. The generation itself was effectively instantaneous: the
slowest of the four stages on a representative deliverable completed in under two
ten thousandths of a second, and the largest artifact produced was a 1000 byte
paper, far below the 200{,}000 byte cap. Three artifacts from the pipeline are
saved in the execution record as concrete evidence of the run, namely the
generated instruction set, the captured execution log, and the auto assembled
paper. The speed of this step is precisely the point developed in the discussion:
producing a well formed candidate is cheap, so the value of the workflow cannot
lie there.

\begin{figure}[H]
    \hspace{-0.8cm}%
    \IfFileExists{Images/acceleration-waterfall.png}
      {\includegraphics[width=1.1\linewidth, trim={0 0 0 0}, clip]{Images/acceleration-waterfall.png}}
      {\figslot{Schedule acceleration waterfall, 30 to 12 days}}
    \caption{Schedule acceleration waterfall from a 30 day baseline to a 12 day automated schedule, a 2.5x reduction.}
    \label{figwaterfall}
\end{figure}

\subsection{The VVUQ gate, the core result}

The central result is the behavior of the gate. Three shipped examples first
exercised each dimension on a clean deliverable. Verification passed all six
structural checks for a fraction of 1.0. Validation reported perfect agreement of
1.0 against the reference with a worst case relative error of 0.033, inside the
0.05 bound, and a recorded human review. Uncertainty, measured across three runs,
reported a maximum coefficient of variation of 0.0068, comfortably inside the
0.10 bound, and the candidate was accepted. The accepted deliverable clears every
spoke of the threshold ring shown in \autoref{figwheel}. The decisive evidence,
however, is the full decision surface in \autoref{tab:results_gate}, where the
gate was exercised across six cases. It accepted exactly one and blocked five,
and one of those five blocks escalated to a human because the runs diverged. The
single most important row is case 6: a candidate that exceeded the 200{,}000 byte
cap failed exactly one of the six verification checks, which dropped the
verification fraction to 0.83, and that single failure was sufficient to block
the deliverable even though everything else about it was well formed. That
asymmetry, a complete and valid looking candidate refused on one failed check out
of six, is the thesis in a single line: the gate is deliberately harder to
satisfy than the generator is to run.

\begin{figure}[H]
    \hspace{-0.8cm}%
    \IfFileExists{Images/vvuq-assurance-wheel.png}
      {\includegraphics[width=1.1\linewidth, trim={0 0 0 0}, clip]{Images/vvuq-assurance-wheel.png}}
      {\figslot{VVUQ assurance wheel, accepted candidate versus threshold ring}}
    \caption{VVUQ assurance wheel: the accepted deliverable clears every spoke, while each blocked case fails the spoke shown.}
    \label{figwheel}
\end{figure}

\begin{table}[H]
\caption{The full VVUQ gate decision surface (1 accept, 5 block, 1 escalate).}
\label{tab:results_gate}
\centering
\begin{tabularx}{\textwidth}{C{0.6cm} L{4.2cm} C{1.3cm} C{1.3cm} >{\raggedright\arraybackslash}X}
\toprule
\textbf{\#} & \textbf{Scenario} & \textbf{Accept} & \textbf{Escal.} & \textbf{Blocking reason} \\
\midrule
  1 & clean, human reviewed, low CV     & yes & no  & none \\
  2 & no human review recorded          & no  & no  & human review not recorded \\
  3 & divergent runs, high CV           & no  & yes & max cv 0.365 above 0.1 \\
  4 & observed disagrees with reference & no  & no  & agreement 0.00; rel err 0.600 \\
  5 & only 2 runs, need 3               & no  & no  & only 2 runs, need 3 \\
  6 & artifact exceeds 200K cap         & no  & no  & verification fraction 0.83 \\
\bottomrule
\end{tabularx}
\end{table}

\subsection{Stage 1 automation}

The Stage 1 capabilities ran around the pipeline and all behaved as designed.
Triple run consensus converged: the acceleration factor metric showed a
coefficient of variation of 0.0007 and the automated days metric a coefficient of
variation of 0.0001 across the three runs, with zero metrics flagged as
divergent. Robust ingestion degraded gracefully: with the network unreachable the
web search client exhausted its four attempt retry budget and then returned a
clearly labeled offline stub rather than raising, and with the PDF backend absent
the processor reported its clean non fatal status and ran only the size based
chunk estimate. Autochunking was verified to be lossless: the chunker is byte
identical on line boundary splits and on ASCII hard splits, both independently
reconstructed and confirmed. One genuine limitation was discovered here and is
carried honestly into \autoref{sec:limitations}: a single oversized line of
multibyte UTF-8 text loses one character at each internal split boundary, six
characters in the tested case, because of a decode that ignores errors after a
byte level cut. The scheduler planned twenty four evenly spaced commit slots at
an hourly interval of 3600 seconds. These outcomes connect the platform to the
prior triplicate and digital twin work that motivated triple run
consensus~\cite{kawchak2025pdactwin}.

\subsection{Stage 2 reference}

The Stage 2 reference was executed as a planning artifact. The 2030 pancreatic
cancer hybrid pilot is expressed as a seven event timeline: a 60 second eight arm
robotic Whipple procedure on day 0, followed by six 28 day adjuvant medicine
cycles that run through day 168, for a total of 168 regimen days, with human
oversight required throughout. The lights off factory controller exercised its
full safety surface correctly. On a clean batch it completed all tasks and
returned to its dark idle state; when cumulative faults exceeded the configured
budget it emergency stopped; and when an interlock was left unsatisfied it blocked
before any task ran. Among the four default interlocks is
\texttt{vvuq\_gate\_online}, so the Stage 1 assurance gate is a precondition for
any Stage 2 batch. This is the concrete mechanism by which the software assurance
of Stage 1 becomes the physical safety of Stage 2, and it is why the gate, not
the generator, is the component that scales across both stages.

\subsection{Per commit processing evidence}

Finally, the execution run was itself recorded one commit per section, and that
record is evidence that real time, per commit assurance is practical at the speed
the pipeline runs. The eight execution commits are listed in
\autoref{tab:results_exec_commits} and drawn as a timeline in
\autoref{figtimeline}. The producing commits land within one to two minutes of
each other, the documentation and the two closing commits take slightly longer,
and the whole run spans roughly eighteen minutes. Because each section was
committed and pushed as soon as it was verified, a regression in any section would
have been attributable to the commit that introduced it rather than discovered in
a single terminal pass, which is the per commit advantage over prior batch
oriented generation.

\begin{table}[H]
\caption{Autonomous execution commits producing the run record.}
\label{tab:results_exec_commits}
\centering
\begin{tabularx}{\textwidth}{C{0.6cm} L{1.9cm} >{\raggedright\arraybackslash}X C{1.7cm} C{2.0cm}}
\toprule
\textbf{\#} & \textbf{Hash} & \textbf{Section} & \textbf{Lines} & \textbf{Time (UTC)} \\
\midrule
  1 & d7721a7 & Foundation (env, tests, lint)  & +387     & 04:20:42 \\
  2 & b146ef7 & Pipeline (five methods)        & +322     & 04:22:43 \\
  3 & c340cef & VVUQ (gate surface)            & +223     & 04:24:05 \\
  4 & f5cb95b & Stage 1 automation             & +238     & 04:26:26 \\
  5 & 2eb5cf0 & Physical-AI Stage 2 (PDAC)     & +176     & 04:27:42 \\
  6 & ddd41e9 & Comprehensive README           & +401 -1  & 04:32:08 \\
  7 & 9437292 & Fix errors (second to last)    & +7 -6    & 04:35:34 \\
  8 & 57aa8db & Repository updates (last)      & +102 -13 & 04:38:44 \\
\bottomrule
\end{tabularx}
\end{table}

\begin{figure}[H]
\centering
\begin{verbatim}
UTC       Commit   Section                        Delta     Gap
--------  -------  -----------------------------  --------  ------
04:20:42  d7721a7  Foundation (env, tests, lint)  +387      start
04:22:43  b146ef7  Pipeline (five methods)        +322      2m01s
04:24:05  c340cef  VVUQ (gate surface)            +223      1m22s
04:26:26  f5cb95b  Stage 1 automation             +238      2m21s
04:27:42  2eb5cf0  Physical-AI Stage 2 (PDAC)     +176      1m16s
04:32:08  ddd41e9  Comprehensive README           +401 -1   4m26s
04:35:34  9437292  Fix errors (second to last)    +7 -6     3m26s
04:38:44  57aa8db  Repository updates (last)      +102 -13  3m10s
--------  -------  -----------------------------  --------  ------
8 commits over 18m02s, all pushed in real time in one pull request
\end{verbatim}
\caption{Execution commit timeline, one section per commit, pushed in real time.}
\label{figtimeline}
\end{figure}

The figures in this section are grounded in the imagegen Python scripts listed in
\autoref{tab:results_figs}; the numeric arrays are taken from the script literals
rather than from any rendered image, and each figure is referenced in the prose
above rather than restated~\cite{sel2025vvuq}.

\begin{table}[H]
\caption{Results figures grounded in the imagegen Python scripts.}
\label{tab:results_figs}
\centering
\begin{tabularx}{\textwidth}{L{4.8cm} >{\raggedright\arraybackslash}X}
\toprule
\textbf{Script directory (under \texttt{imagegen/})} & \textbf{Renders as} \\
\midrule
  \texttt{02-acceleration-waterfall} & A schedule waterfall, 30 to 12 days via reductions of 10, 5, and 3, with the audit formula. \\
  \texttt{04-vvuq-assurance-wheel} & A radar wheel of the six normalized assurance spokes against the threshold ring. \\
  \texttt{06-test-coverage-treemap} & A treemap of 51 of 51 tests across the 8 modules, areas 17, 8, 5, 5, 4, 4, 4, 4. \\
\bottomrule
\end{tabularx}
\end{table}
